\setcounter{section}{6}






  \zwischenueberschrift{Exaktheit}


\inputdefinition
{}
{





Es sei $X$ ein
\definitionsverweis {topologischer Raum}{}{,}
es seien ${ \mathcal  F   }_n$
\definitionsverweis {Garben}{}{}
von
\definitionsverweis {kommutativen Gruppen}{}{}
auf $X$ und es seien
\maabb {\varphi_n} { { \mathcal  F   }_{n-1} } { { \mathcal  F   }_n
} {}
\definitionsverweis {Homomorphismen}{}{.}
Man sagt, dass ein
\definitionswort {Garbenkomplex}{}
vorliegt, wenn

\mavergleichskettedisp
{\vergleichskette
{ \operatorname{bild} \varphi_{n}
}
{ \subseteq} { \operatorname{kern} \varphi_{n+1}
}
{ } {
}
{ } {
}
{ } {
}
}
{}{}{}
gilt.





}


\inputdefinition
{}
{





Es sei $X$ ein
\definitionsverweis {topologischer Raum}{}{}
und es sei ${ \mathcal  F   }_\bullet$ ein
\definitionsverweis {Komplex}{}{}
von
\definitionsverweis {Garben}{}{}
von
\definitionsverweis {kommutativen Gruppen}{}{}
auf $X$. Man sagt, dass der Komplex
\definitionswort {exakt}{}
ist, wenn

\mavergleichskettedisp
{\vergleichskette
{ \operatorname{bild} \varphi_{n}
}
{ =} { \operatorname{kern} \varphi_{n+1}
}
{ } {
}
{ } {
}
{ } {
}
}
{}{}{}
für alle

\mavergleichskette
{\vergleichskette
{n
}
{ \in }{\Z
}
{  }{
}
{    }{
}
{   }{
}
}
{}{}{}
gilt.





}





\inputfaktbeweis
{Garbe/Gruppen/Exaktheit/Charakterisierungen/Fakt}
{Lemma}
{}
{



\faktsituation {Es sei $X$ ein
\definitionsverweis {topologischer Raum}{}{}
und es sei

\mathdisp {{ \mathcal  F   }  \stackrel{  }{\longrightarrow }  { \mathcal  G   }  \stackrel{   }{\longrightarrow }   { \mathcal  H   }} {  }

ein
\definitionsverweis {Komplex}{}{}
von
\definitionsverweis {Garben von kommutativen Gruppen}{}{}
auf $X$.}
\faktfolgerung {Dann ist der Komplex genau dann
\definitionsverweis {exakt}{}{,}
wenn für jeden Punkt

\mavergleichskette
{\vergleichskette
{P
}
{ \in }{ X
}
{  }{
}
{    }{
}
{   }{
}
}
{}{}{}
der Komplex

\mathdisp {{ \mathcal  F   }_P  \stackrel{  }{\longrightarrow }  { \mathcal  G   }_P  \stackrel{   }{\longrightarrow }   { \mathcal  H   }_P} {  }

\definitionsverweis {exakt}{}{}
ist.}
\faktzusatz {}
\faktzusatz {}





}
{





Wir benennen die Situation mit

\mathdisp {{ \mathcal  F   }  \stackrel{ \alpha }{\longrightarrow }  { \mathcal  G   }  \stackrel{  \beta }{\longrightarrow }   { \mathcal  H   }} { . }

Nach

Korollar 4.11

liegt ein
\definitionsverweis {Garbenkomplex}{}{}
genau dann vor, wenn sämtliche Halmabbildungen
\definitionsverweis {Komplexe}{}{}
sind. Es sei der Komplex exakt, also

\mavergleichskettedisp
{\vergleichskette
{ \operatorname{bild} \alpha
}
{ =} { \operatorname{kern} \beta
}
{ } {
}
{ } {
}
{ } {
}
}
{}{}{.}
Es sei

\mavergleichskette
{\vergleichskette
{P
}
{ \in }{ X
}
{  }{
}
{    }{
}
{   }{
}
}
{}{}{}
fixiert und sei

\mavergleichskette
{\vergleichskette
{s
}
{ \in }{ { \mathcal  G   }_P
}
{  }{
}
{    }{
}
{   }{
}
}
{}{}{}
mit

\mavergleichskette
{\vergleichskette
{ \beta_P(s)
}
{ = }{ 0
}
{  }{
}
{    }{
}
{   }{
}
}
{}{}{.}
Dann gibt es eine offene Umgebung $U$ von $P$ auf der $s$ durch einen Schnitt $s$ repräsentiert wird und eine kleinere offene Umgebung

\mavergleichskette
{\vergleichskette
{V
}
{ \subseteq }{ U
}
{  }{
}
{    }{
}
{   }{
}
}
{}{}{,}
worauf

\mavergleichskette
{\vergleichskette
{ \beta_V(s \vert_V )
}
{ = }{ 0
}
{  }{
}
{    }{
}
{   }{
}
}
{}{}{}
ist. Das Element
\zusatzklammer {wir nennen die Einschränkung wieder $s$} {} {}

\mavergleichskette
{\vergleichskette
{s
}
{ \in }{ { \mathcal  G   }  \left(   V   \right)
}
{  }{
}
{    }{
}
{   }{
}
}
{}{}{}
gehört also zum Kern von $\beta_V$ und daher zum
\zusatzklammer {Garben} {-} {}Bild von $\alpha$. D.h. es gibt eine offene Umgebung

\mavergleichskette
{\vergleichskette
{P
}
{ \in }{ W
}
{ \subseteq }{ V
}
{    }{
}
{   }{
}
}
{}{}{,}
auf der $s$ im Bild von
\maabbdisp {\alpha_W} { { \mathcal  F   }  \left(   W   \right)  } { { \mathcal  G   }  \left(   W   \right)
} {}
liegt. Daher liegt auch der Keim $s$ im Bild von $\alpha_P$.




}





\inputdefinition
{}
{





Ein
\definitionsverweis {exakter}{}{}
\definitionsverweis {Komplex}{}{}

\mathdisp {0 \, \stackrel{  }{\longrightarrow} \,  { \mathcal  F   }   \, \stackrel{  }{ \longrightarrow}   \,  { \mathcal  G   }   \, \stackrel{  }{\longrightarrow} \,  { \mathcal  H   }  \,\stackrel{  } {\longrightarrow}  \, 0} {  }

von
\definitionsverweis {Garben von kommutativen Gruppen}{}{}
auf einem
\definitionsverweis {topologischen Raum}{}{}
$X$ heißt
\definitionswort {kurze exakte Sequenz}{.}





}
Hierbei ist insbesondere die vordere Abbildung injektiv und die hintere Abbildung
\zusatzklammer {Garben} {} {-}surjektiv.





\inputfaktbeweis
{Topologische Gruppen/Kommutativ/Kurze exakte Sequenz/Lokal stetiger Schnitt/Garbensequenz/Fakt}
{Lemma}
{}
{



\faktsituation {Es sei

\mathdisp {0 \, \stackrel{  }{\longrightarrow} \,  F   \, \stackrel{  }{ \longrightarrow}   \,  G   \, \stackrel{  }{\longrightarrow} \,  H \,\stackrel{  } {\longrightarrow}  \, 0} {  }

eine
\definitionsverweis {kurze exakte Sequenz}{}{}
von kommutativen
\definitionsverweis {topologischen Gruppen}{}{}
\zusatzklammer {mit stetigen Gruppenhomomorphismen} {} {.}}
\faktvoraussetzung {Es trage $F$ die
\definitionsverweis {induzierte Topologie}{}{}
von $G$ und die Surjektion
\maabb {p} { G } { H
} {}
habe die Eigenschaft, dass es zu jedem Element

\mavergleichskette
{\vergleichskette
{h
}
{ \in }{ H
}
{  }{
}
{    }{
}
{   }{
}
}
{}{}{}
eine offene Umgebung

\mavergleichskette
{\vergleichskette
{h
}
{ \in }{ W
}
{ \subseteq }{ H
}
{    }{
}
{   }{
}
}
{}{}{}
und einen
\definitionsverweis {stetigen Schnitt}{}{}
zu $p$ gibt.}
\faktfolgerung {Dann ist für jeden
\definitionsverweis {topologischen Raum}{}{}
$X$ die zugehörige Sequenz der Garben der stetigen Abbildungen in diese Gruppen, also

\mathdisp {0 \, \stackrel{  }{\longrightarrow} \, C^0(-,F)  \, \stackrel{  }{ \longrightarrow}   \, C^0(-,G)    \, \stackrel{  }{\longrightarrow} \,  C^0(-,H)   \,\stackrel{  } {\longrightarrow}  \, 0} { , }

ebenfalls
\definitionsverweis {exakt}{}{.}}
\faktzusatz {}
\faktzusatz {}





}
{





Es ist klar, dass ein
\definitionsverweis {Komplex}{}{}
von
\definitionsverweis {Garben von kommutativen Gruppen}{}{}
auf $X$ vorliegt. Die Injektivität links ist ebenfalls klar. Zur Exaktheit in der Mitte: Wenn zu einer offenen Menge

\mavergleichskette
{\vergleichskette
{U
}
{ \subseteq }{ X
}
{  }{
}
{    }{
}
{   }{
}
}
{}{}{}
eine stetige Abbildung
\maabb {\varphi} { U } { G
} {}
die Eigenschaft besitzt, dass
\mathl{p \circ \varphi}{} die Nullabbildung ist, so liegt das Bild von $\varphi$ in $F$. Da $F$ die
\definitionsverweis {induzierte Topologie}{}{}
von $G$ trägt, ist auch die Abbildung
\maabb {\varphi} { U } { F
} {}
stetig. Zur Garbensurjektivität rechts: Es sei

\mavergleichskette
{\vergleichskette
{P
}
{ \in }{ X
}
{  }{
}
{    }{
}
{   }{
}
}
{}{}{}
ein Punkt und
\maabbdisp {\psi} { V} {H
} {}
eine auf einer offenen Umgebung von $P$ definierte stetige Abbildung nach $H$. Es sei

\mavergleichskette
{\vergleichskette
{ \psi(P)
}
{ = }{ h
}
{  }{
}
{    }{
}
{   }{
}
}
{}{}{.}
Nach Voraussetzung gibt es eine offene Umgebung

\mavergleichskette
{\vergleichskette
{h
}
{ \in }{ W
}
{ \subseteq }{ H
}
{    }{
}
{   }{
}
}
{}{}{}
und einen Schnitt
\maabb {s} { W} {G
} {}
mit

\mavergleichskette
{\vergleichskette
{ p \circ s
}
{ = }{
\operatorname{Id}_{  W  }
}
{  }{
}
{    }{
}
{   }{
}
}
{}{}{.}
Wir betrachten

\mavergleichskettedisp
{\vergleichskette
{U
}
{ \defeq} { V \cap \psi^{-1}(W)
}
{ } {
}
{ } {
}
{ } {
}
}
{}{}{.}
Dann ist
\mathl{s \circ \psi}{}
\zusatzklammer {eingeschränkt auf $U$} {} {}
ein stetiger Schnitt von $G$, der unter $p$ auf $\psi$ abgebildet wird.




}




\inputbeispiel{}
{





Wir betrachten die kurze exakte \stichwort {Exponentialsequenz} {}

\mathdisp {0 \, \stackrel{  }{\longrightarrow} \,  2 \pi { \mathrm i}\Z  \, \stackrel{  }{ \longrightarrow}   \,  {\mathbb C}    \, \stackrel{ \operatorname{exp} }{\longrightarrow} \,  {\mathbb C} ^{\times}  \,\stackrel{  } {\longrightarrow}  \, 0} {  }

von
\definitionsverweis {topologischen Gruppen}{}{.}
Die Exaktheit in der Mitte beruht auf

Satz 21.5 (Analysis (Osnabrück 2021-2023))  (2)
,
die Homomorphieeigenschaft beruht auf
der
Funktionalgleichung der Exponentialfunktion

und die
\definitionsverweis {komplexe Exponentialfunktion}{}{}
bildet nach

Satz 21.6 (Analysis (Osnabrück 2021-2023))

surjektiv auf ${\mathbb C} \setminus \{0\}$ ab
\zusatzklammer {sie ist eine
\definitionsverweis {Überlagerung}{}{,}
siehe

Beispiel 21.3 (Funktionentheorie (Osnabrück 2023-2024))
} {} {.}
Da es lokal einen Logarithmus gibt, sind die Voraussetzungen von

Lemma 6.5

erfüllt. Somit gibt es zu jedem topologischen Raum $X$ eine
\definitionsverweis {kurze exakte Garbensequenz}{}{}

\mathdisp {0 \, \stackrel{  }{\longrightarrow} \, C^0(-, \Z)  \, \stackrel{  }{ \longrightarrow}   \, C^0(-, {\mathbb C} )    \, \stackrel{  }{\longrightarrow} \,  C^0(-, {\mathbb C} ^{\times} )   \,\stackrel{  } {\longrightarrow}  \, 0} { , }

die die
\zusatzklammer {stetige komplexe} {} {}
\stichwort {Exponentialsequenz} {} heißt. Links steht die lokal konstante Garbe mit Werten in $\Z$, in der Mitte die Garbe der komplexwertigen stetigen Funktionen und rechts die Garbe der nullstellenfreien komplexwertigen stetigen Funktionen. Wenn man

\mavergleichskette
{\vergleichskette
{X
}
{ = }{ {\mathbb C} ^{\times}
}
{  }{
}
{    }{
}
{   }{
}
}
{}{}{}
setzt, so ist die globale Auswertung der hinteren Abbildung nicht surjektiv, da die Identität nicht im Bild liegt.





}





  \zwischenueberschrift{Globale Auswertung}





\inputfaktbeweis
{Garbe/Kommutative Gruppen/Komplex/Globale Auswertung/Fakt}
{Lemma}
{}
{



\faktsituation {Es sei $X$ ein
\definitionsverweis {topologischer Raum}{}{}}
\faktvoraussetzung {und sei

\mathdisp {{ \mathcal  F   } \stackrel{d}{\longrightarrow } { \mathcal  G   } \stackrel{d'}{\longrightarrow} { \mathcal  H   }} {  }

ein
\definitionsverweis {Komplex}{}{}
von
\definitionsverweis {Garbenhomomorphismen}{}{}
von
\definitionsverweis {Garben}{}{}
von
\definitionsverweis {kommutativen Gruppen}{}{}
auf $X$.}
\faktfolgerung {Dann ist auch

\mathdisp {\Gamma   \left(   X ,  { \mathcal  F   }   \right) \longrightarrow \Gamma   \left(   X ,  { \mathcal  G   }   \right) \longrightarrow \Gamma   \left(   X ,  { \mathcal  H   }   \right)} {  }

ein Komplex.}
\faktzusatz {}
\faktzusatz {}





}
{





Die Voraussetzung bedeutet einfach, dass
\mathl{d' \circ d}{} die Nullabbildung ist. Dann ist insbesondere die globale Auswertung die Nullabbildung.




}






\inputfaktbeweis
{Garbe/Kommutative Gruppen/Linksexakt/Fakt}
{Lemma}
{}
{



\faktsituation {Es sei $X$ ein
\definitionsverweis {topologischer Raum}{}{.}}
\faktvoraussetzung {Es sei

\mathdisp {0 \longrightarrow { \mathcal  F   } \stackrel{d}{ \longrightarrow } { \mathcal  G   } \stackrel{d'}{\longrightarrow } { \mathcal  H   }} {  }

ein exakter Komplex von
\definitionsverweis {Garbenhomomorphismen}{}{}
von
\definitionsverweis {Garben}{}{}
von
\definitionsverweis {kommutativen Gruppen}{}{}
auf $X$.}
\faktfolgerung {Dann ist auch der Komplex

\mathdisp {0 \longrightarrow \Gamma   \left(   X ,  { \mathcal  F   }   \right) \longrightarrow \Gamma   \left(   X ,  { \mathcal  G   }   \right) \longrightarrow \Gamma   \left(   X ,  { \mathcal  H   }   \right)} {  }

exakt.}
\faktzusatz {}
\faktzusatz {}





}
{





Dass ein Komplex vorliegt ist klar nach

Lemma 6.7
.
Die Exaktheit bedeutet, dass für jeden Punkt

\mavergleichskette
{\vergleichskette
{P
}
{ \in }{ X
}
{  }{
}
{    }{
}
{   }{
}
}
{}{}{}
der Komplex

\mathdisp {0 \longrightarrow { \mathcal  F   }_P \longrightarrow { \mathcal  G   }_P \longrightarrow { \mathcal  H   }_P} {  }

der Halme exakt ist. Sei

\mavergleichskette
{\vergleichskette
{s
}
{ \in }{ \Gamma   \left(   X ,  { \mathcal  F   }   \right)
}
{  }{
}
{    }{
}
{   }{
}
}
{}{}{}
und

\mavergleichskette
{\vergleichskette
{ d(s)
}
{ = }{ 0
}
{  }{
}
{    }{
}
{   }{
}
}
{}{}{}
in
\mathl{\Gamma   \left(   X ,  { \mathcal  G   }   \right)}{.} Dann ist

\mavergleichskette
{\vergleichskette
{ d(s)_P
}
{ = }{ 0
}
{  }{
}
{    }{
}
{   }{
}
}
{}{}{}
in jedem Punkt und somit ist

\mavergleichskette
{\vergleichskette
{ s_P
}
{ = }{ 0
}
{  }{
}
{    }{
}
{   }{
}
}
{}{}{}
für jeden Punkt. Also ist

\mavergleichskette
{\vergleichskette
{ s
}
{ = }{ 0
}
{  }{
}
{    }{
}
{   }{
}
}
{}{}{}
nach

Lemma 4.4

und die linke Abbildung ist injektiv. Es sei nun

\mavergleichskette
{\vergleichskette
{ t
}
{ \in }{ \Gamma   \left(   X ,  { \mathcal  G   }   \right)
}
{  }{
}
{    }{
}
{   }{
}
}
{}{}{}
mit

\mavergleichskette
{\vergleichskette
{ d'(t)
}
{ = }{ 0
}
{  }{
}
{    }{
}
{   }{
}
}
{}{}{}
in
\mathl{\Gamma   \left(   X ,  { \mathcal  H   }   \right)}{.} Die Exaktheit in den Halmen bedeutet, dass für jeden Punkt $P$ der Keim $t_P$ zu
\mathl{{ \mathcal  F   }_P}{} gehört. Daraus folgt mit

Aufgabe 5.5
,
dass $t$ selbst zu ${ \mathcal  F   }$ gehört.




}



Die vorstehende Aussage bedeutet, dass die globale Auswertung einer Garbe von abelschen Gruppen ein
\zusatzklammer {kovarianter, additiver} {} {}
\definitionsverweis {linksexakter Funktor}{}{}
ist.





  \zwischenueberschrift{Rückzug und Vorschub}

Bisher haben wir nur Garben und ihre Beziehungen untereinander auf einem gegebenen topologischen Raum behandelt. Wir betrachten nun den Fall, wo topologische Räume durch eine stetige Abbildung miteinander verbunden sind.


\inputdefinition
{}
{





Zu einer
\definitionsverweis {stetigen Abbildung}{}{}
\maabbdisp {\varphi} { X } { Y
} {}
und einer
\definitionsverweis {Prägarbe}{}{}
${ \mathcal  F   }$ auf $X$ nennt man die durch

\mavergleichskettedisp
{\vergleichskette
{ \left(  \varphi_*  \right) { \mathcal  F   }  (U)
}
{ \defeq} { { \mathcal  F   } \left(  \varphi^{-1}(U)  \right)
}
{ } {
}
{ } {
}
{ } {
}
}
{}{}{}
gegebene Prägarbe auf $Y$ die unter $\varphi$
\definitionswort {vorgeschobene Prägarbe}{.}





}
Da zu offenen Mengen

\mavergleichskette
{\vergleichskette
{V
}
{ \subseteq }{W
}
{  }{
}
{    }{
}
{   }{
}
}
{}{}{}
auch

\mavergleichskette
{\vergleichskette
{\varphi^{-1}(V)
}
{ \subseteq }{\varphi^{-1}(W)
}
{  }{
}
{    }{
}
{   }{
}
}
{}{}{}
gilt, hat man natürliche Restriktionsabbildungen und erhält somit in der Tat eine Prägarbe.





\inputfaktbeweis
{Stetige Abbildung/Garbe/Vorschub/Fakt}
{Lemma}
{}
{



\faktsituation {Zu einer
\definitionsverweis {stetigen Abbildung}{}{}
\maabbdisp {\varphi} { X } { Y
} {}
und einer
\definitionsverweis {Garbe}{}{}
${ \mathcal  F   }$ auf $X$ ist die
\definitionsverweis {vorgeschobene Prägarbe}{}{}

\mathl{\varphi_* { \mathcal  F   }}{}}
\faktfolgerung {eine
\definitionsverweis {Garbe}{}{.}}
\faktzusatz {}
\faktzusatz {}





}
{





Es sei

\mavergleichskettedisp
{\vergleichskette
{V
}
{ =} { \bigcup_{i \in I} V_i
}
{ } {
}
{ } {
}
{ } {
}
}
{}{}{}
eine offene Überdeckung einer offenen Menge

\mavergleichskette
{\vergleichskette
{V
}
{ \subseteq }{ Y
}
{  }{
}
{    }{
}
{   }{
}
}
{}{}{.}
Dann bilden die

\mathbed {\varphi^{-1} (V_i)} {}
 {i \in I} {}
 {} {} {} {,}
eine offene Überdeckung von $\varphi^{-1}(V)$. Es seien
\mathl{s,t \in \varphi_*{ \mathcal F  }(V)}{} mit

\mavergleichskette
{\vergleichskette
{ s \vert_{V_i \cap V_j }
}
{ = }{ t \vert_{V_i \cap V_j }
}
{  }{
}
{    }{
}
{   }{
}
}
{}{}{}
gegeben. Dies bedeutet unmittelbar
\mathl{s,t \in { \mathcal F  } \left(  \varphi^{-1}(V)  \right)}{} und

\mavergleichskettedisp
{\vergleichskette
{ s\vert_{  \varphi^{-1}(V_i) \cap  \varphi^{-1}(V_j)  }
}
{ =} { s\vert_{  \varphi^{-1}(V_i \cap V_j)  }
}
{ =} { t\vert_{  \varphi^{-1}(V_i \cap V_j)  }
}
{ =} { t\vert_{  \varphi^{-1}(V_i) \cap  \varphi^{-1}(V_j)  }
}
{ } {
}
}
{}{}{.}
Daher ist
\zusatzklammer {nach der ersten Garbeneigenschaft von ${ \mathcal  F   }$} {} {}

\mavergleichskette
{\vergleichskette
{s
}
{ = }{t
}
{  }{
}
{    }{
}
{   }{
}
}
{}{}{}
in ${ \mathcal  F   } \left(  \varphi^{-1} (V)   \right)$, also

\mavergleichskette
{\vergleichskette
{s
}
{ = }{t
}
{  }{
}
{    }{
}
{   }{
}
}
{}{}{}
in $\varphi_* { \mathcal  F   }  (V)$.

Es seien nun
\mathl{s_i \in { \mathcal  F   }  (V_i)}{} mit

\mavergleichskettedisp
{\vergleichskette
{ s_i \vert_{V_i \cap V_j }
}
{ =} { s_j\vert_{V_i \cap V_j }
}
{ } {
}
{ } {
}
{ } {
}
}
{}{}{.}
Dies bedeutet zurückübersetzt nach $X$ unmittelbar, dass kompatible Schnitte in
\mathl{{ \mathcal  F   } \left(  \varphi^{-1}(V_i)   \right)}{} vorliegen, denen ein Schnitt in

\mavergleichskette
{\vergleichskette
{ { \mathcal  F   } \left(  \varphi^{-1}(V)   \right)
}
{ = }{ \varphi_* { \mathcal  F   }  (V)
}
{  }{
}
{    }{
}
{   }{
}
}
{}{}{}
entspricht.




}






\inputfaktbeweis
{Stetige Abbildung/Prägarbe/Vorschub/Halme/Fakt}
{Lemma}
{}
{



\faktsituation {Zu einer
\definitionsverweis {stetigen Abbildung}{}{}
\maabbdisp {\varphi} { X } { Y
} {,}
einem Punkt

\mavergleichskette
{\vergleichskette
{Q
}
{ \in }{ Y
}
{  }{
}
{    }{
}
{   }{
}
}
{}{}{}
und einer
\definitionsverweis {Prägarbe}{}{}
${ \mathcal  F   }$ auf $X$ ist der
\definitionsverweis {Halm}{}{}
der
\definitionsverweis {vorgeschobenen Prägarbe}{}{}

\mathl{\varphi_* { \mathcal  F   }}{} im Punkt $Q$ gleich}
\faktfolgerung {
\mavergleichskettedisphandlinks
{\vergleichskettedisphandlinks
{ \operatorname{colim}\,_{Q \in V } { \mathcal  F   } \left(  \varphi^{-1}(V)  \right)
}
{ =} { \operatorname{colim}\,_{ \left\{  U \subseteq X  \mid  \text{ es gibt eine offene Umgebung } Q \in V \text{ mit } \varphi^{-1}(V) \subseteq U         \right\}  } { \mathcal  F   } \left(  U   \right)
}
{ } {
}
{ } {
}
{ } {
}
}
{}{}{.}}
\faktzusatz {}
\faktzusatz {}




 }
{ Siehe
Aufgabe 6.5
. }


Der Halm der vorgeschobenen Prägarbe ist also der Halm der Ausgangsgarbe in einem Filter
\zusatzklammer {nämlich dem Urbildfilter des Umgebungsfilters
\mathl{U(Q)}{}} {} {,}
aber im Allgemeinen nicht in einem Punkt.


\inputdefinition
{}
{





Zu einer
\definitionsverweis {stetigen Abbildung}{}{}
\maabbdisp {\varphi} { X } { Y
} {}
und einer
\definitionsverweis {Prägarbe}{}{}
${ \mathcal  G   }$ auf $Y$ nennt man auf einer offenen Menge

\mavergleichskette
{\vergleichskette
{U
}
{ \subseteq }{ X
}
{  }{
}
{    }{
}
{   }{
}
}
{}{}{}
durch

\mathdisp {\operatorname{colim}_{  { \mathcal  G   }  \left(   V   \right)  }\, V \subseteq Y,\, U \subseteq \varphi^{-1} (V)} {  }

gegebene Prägarbe auf $X$ die unter $\varphi$
\definitionswort {zurückgezogene Prägarbe}{.}





}


\inputdefinition
{}
{





Zu einer
\definitionsverweis {stetigen Abbildung}{}{}
\maabb {\varphi} { X } { Y
} {}
und einer
\definitionsverweis {Garbe}{}{}
${ \mathcal  G   }$ auf $Y$ nennt man die
\definitionsverweis {Vergarbung}{}{}
der
\definitionsverweis {zurückgezogenen Prägarbe}{}{}
die
\definitionswort {zurückgezogene Garbe}{.}





}
Sie wird mit $\varphi^{-1} { \mathcal G  }$ bezeichnet.





\inputfaktbeweis
{Stetige Abbildung/Garbe/Rückzug/Halme/Fakt}
{Lemma}
{}
{



\faktsituation {Zu einer
\definitionsverweis {stetigen Abbildung}{}{}
\maabb {\varphi} { X } { Y
} {}
und einer
\definitionsverweis {Garbe}{}{}
${ \mathcal  G   }$ auf $Y$}
\faktfolgerung {ist der
\definitionsverweis {Halm}{}{}
der
\definitionsverweis {zurückgezogenen Garbe}{}{}
in einem Punkt

\mavergleichskette
{\vergleichskette
{P
}
{ \in }{ X
}
{  }{
}
{    }{
}
{   }{
}
}
{}{}{}
gleich dem Halm von ${ \mathcal  G   }$ in
\mathl{\varphi(P)}{.}}
\faktzusatz {}
\faktzusatz {}




 }
{ Siehe
Aufgabe 6.6
. }
